% Offline descriptive follow-up, completed 2026-09-30.
A result-informed loss sensitivity analysis compared the fixed two-request checks without selecting a new deployment rule. When total request loss is omitted ($K_p$=0), hydraulic Same and View have empirical losses (4+17c)/54 and (5+10c)/54. Their difference, View minus Same, is (1-7c)/54 and favors View only when c>1/7: its six additional accepted correct diagnoses accompany one additional accepted error. Compared with accepting every F1 output, the empirical break-even referral losses are 6/17 for Same and 1/2 for View. Neither check improves every endpoint.

In the bearing study, the corresponding losses are (58+6c)/72 and (21+42c)/72. View has lower observed loss than Same throughout 0$\leq$c$\leq$1, but this comparison alone does not establish useful deployment: rejecting every diagnosis has lower observed loss than View when c<7/10, and accepting every F1 diagnosis has lower observed loss when c>20/21. Only between these empirical intersections does View beat both reference actions under this loss specification with $K_p$=0. The retained error fraction remains 21/30 (70.0\%). These are descriptive sample-dependent tradeoffs, not optimized, validated or calibrated industrial decisions.

The decision counts contain two rounds on 27 hydraulic cycles and 36 bearing recordings, not 54 or 72 independent devices. Source and round breakdowns retain the dependence structure. The hydraulic data concern one rig; its 27 condition blocks are not 27 physical rigs, and no rig-level confidence interval is estimated. For Paderborn, class-stratified paired resampling keeps all four regimes and both rounds together for each of nine physical bearings. Its 2.5th--97.5th percentiles describe sensitivity to empirical specimen weights conditional on the observed requests and fixed backgrounds. Shared batch responses leave dependence across bearings, so these bands are not population confidence intervals or evidence of cross-rig generalization.

Round and deletion sensitivity limits any stability claim. Hydraulic Same and View have equal aggregate counts in round 1; the entire six-correct/one-error increment occurs in round 2. Deleting one hydraulic panel moves the View--Same crossing among 0, 1, 1/7; therefore the pooled crossing is not a stable engineering constant.

For bearings, View reduces accepted errors in both rounds (by 20 and 17) but retains -1 and +2 correct diagnoses relative to Same. Its round-1 loss becomes higher than Same for c>20/21; round 2 favors View over the entire displayed range. Deleting each physical bearing still reduces accepted errors (by 29--40), but deleting KI04 reverses the high-cost comparison for c>35/37. Thus error-release reduction survives these particular deletions, whereas a uniformly lower loss does not.

At c=1 the conditional bearing reweighting percentiles for View-minus-Same loss are [-0.0972, 0.0556], spanning both signs. This is a sensitivity to the weights assigned to the nine observed bearings, not a population confidence interval or a validation on additional specimens.
